\ulohahead{programování (C\#)}{Ovladač zařízení programovaným V/V (PIO)}
\begin{zadani}
Pro řadič disku s porty STATUS, COMMAND a DATA napište programem řízenou obsluhu (PIO): vydejte příkaz čtení a přečtěte 512 bajtu polling-em stavového registru. Vysvětlete variantu s přerušením. \emph{(Typ úlohy: 2024-jaro Q2 - ovladač pro řadič disku; 2021-podzim - myš.)}
\end{zadani}

\begin{reseni}
Řadič má registry adresované jako V/V porty; CPU s ním komunikuje instrukcemi in/out (zde abstraktně \texttt{Inb}/\texttt{Outb}).
\begin{lstlisting}
const int DATA = 0x1F0, STATUS = 0x1F7, COMMAND = 0x1F7;
const int BSY = 0x80, DRQ = 0x08;       // bity stavoveho registru

static byte Inb(int port) { /* HW: precti bajt z portu */ return 0; }
static void Outb(int port, byte v) { /* HW: zapis bajt na port */ }

static byte[] ReadSector() {            // PIO ctene polling-em
    Outb(COMMAND, 0x20);                // prikaz READ SECTOR
    var buf = new byte[512];
    for (int i = 0; i < 512; i++) {
        while ((Inb(STATUS) & BSY) != 0) { }   // cekej, az radic neni zaneprazdnen
        while ((Inb(STATUS) & DRQ) == 0) { }   // cekej, az jsou data pripravena
        buf[i] = Inb(DATA);                     // precti bajt dat
    }
    return buf;
}
\end{lstlisting}
\textbf{Polling vs přerušení:} polling \emph{aktivně} cyklí na stavovém registru a pálí CPU. \emph{Přerušení}: řadič po dokončení vyvolá IRQ; procesor přeruší běžící kód, skočí do obslužné rutiny (ISR), ta přečte data a oznámí dokončení (nastaví událost / probudí čekající vlákno). Procesor: dokončí instrukci, uloží kontext, podle vektoru přerušení skočí do ISR; OS: naplánuje obsluhu, po \texttt{iret} obnoví kontext.
\end{reseni}

\begin{jadro}
Registry řadiče = V/V porty (\texttt{in}/\texttt{out}). PIO polling: vydej příkaz, opakovaně čti STATUS (BSY/DRQ), pak čti DATA. Přerušení místo pollingu: IRQ $\to$ ISR přečte data $\to$ probudí čekajícího (nepálí CPU).
\end{jadro}

\past{Polling plýtvá CPU - u pomalých zařízení se použije přerušení. Pořadí bitu: čekej, až BSY=0 a DRQ=1, teprve pak čti DATA. Příkaz se zapisuje do COMMAND, data se čtou z DATA.}
